\textbf{Task.}
Give a Lean-friendly proof plan for a finite-dimensional, endpoint-unified
Riesz-Thorin theorem for complex matrices.

\textbf{Mathematical setting.}
Let \(A : \mathbb C^n \to \mathbb C^m\) be a finite matrix.  For
\(p \in [1,\infty]\), let \(\|x\|_p\) be the finite-dimensional complex
\(\ell^p\) norm, and let
\[
  \|A\|_p = \sup_{x \ne 0} \frac{\|Ax\|_p}{\|x\|_p}.
\]
For \(p_0,p_1,r \in [1,\infty]\) and \(0 \le \theta \le 1\), assume
\[
  \frac1r = (1-\theta)\frac1{p_0}+\theta\frac1{p_1}
\]
with the usual convention \(1/\infty=0\).  Suppose
\[
  \|Ax\|_{p_0}\le M_0\|x\|_{p_0}, \qquad
  \|Ax\|_{p_1}\le M_1\|x\|_{p_1}
\]
for all \(x\), with \(M_0,M_1\ge0\).  The desired conclusion is
\[
  \|Ax\|_r \le M_0^{1-\theta}M_1^\theta\|x\|_r
\]
and hence \(\|A\|_r \le M_0^{1-\theta}M_1^\theta\), including endpoint
cases \(p_i=1,\infty\), \(r=1,\infty\), \(\theta=0,1\).

\textbf{What is already locally proved.}
The Lean development already has:
\begin{itemize}
  \item finite-real strict/interior Riesz-Thorin via Hadamard three-lines for
        Holder-conjugate exponents \(1<p_i,r<\infty\);
  \item endpoint reductions showing that \(\theta=0\) gives \(r=p_0\),
        \(\theta=1\) gives \(r=p_1\), strict-interior \(r=1\) forces
        \(p_0=p_1=1\), and strict-interior \(r=\infty\) forces
        \(p_0=p_1=\infty\);
  \item finite-real wrappers from interpolation data to matrix-bound and
        least-value statements;
  \item special source theorems for the \(1/\infty\) and \(1/2\) endpoint
        rows from Higham equations (6.19)-(6.20).
\end{itemize}

\textbf{Question.}
What is the cleanest mathematically rigorous case split for the fully
endpoint-unified theorem over \(p_0,p_1,r\in[1,\infty]\)?  In particular:
\begin{enumerate}
  \item In the strict case \(0<\theta<1\) and \(1<r<\infty\), prove that any
        endpoint \(p_0\) or \(p_1\) may still be handled by the three-lines
        proof, or explain precisely which additional endpoint boundary lemma is
        needed.
  \item If \(p_0=1\) or \(p_1=\infty\) but \(1<r<\infty\), what dual exponent
        and analytic family should be used so the proof remains valid without
        an \(L^\infty\) dual norming vector?
  \item Give a minimal theorem-dependency ordering suitable for Lean: algebraic
        reciprocal-exponent lemmas first, endpoint-only reductions second,
        one-sided endpoint Hadamard boundary lemmas third, and final wrapper.
  \item Point out hidden assumptions around zero dimensions, zero matrix bounds,
        and real powers \(0^0\).
\end{enumerate}

Please answer as a rigorous proof plan with exact cases, not as prose
encouragement.  Do not assume unproved compactness or norm-attainment facts
unless they are explicitly stated as a dependency.
